% ===========================================================================
% 格点QCD GPU计算管线的物理分析报告
% Physical Analysis Report of the Lattice QCD GPU Pipeline
% ===========================================================================
% Ensemble: beta6.20_mu-0.2770_ms-0.2400_L24x72
% Date: 2026-08-01
% ===========================================================================
\documentclass[11pt,a4paper]{article}

% ── 中文支持 ──
\usepackage[UTF8]{ctex}
\setCJKmainfont{AR PL SungtiL GB}[BoldFont=AR PL UMing CN]
\setCJKsansfont{AR PL KaitiM GB}[BoldFont=AR PL UMing CN]
\setCJKmonofont{AR PL UMing CN}

% ── 数学 ──
\usepackage{amsmath,amssymb,amsfonts}
\usepackage{physics}
\usepackage{braket}
\usepackage{bm}

% ── 排版 ──
\usepackage[margin=2.5cm]{geometry}
\usepackage{booktabs}
\usepackage{graphicx}
\usepackage[colorlinks=true,linkcolor=blue,citecolor=blue,urlcolor=blue]{hyperref}
\usepackage{enumitem}
\usepackage{caption}
\usepackage{subcaption}
\usepackage{multirow}
\usepackage{array}

% ── 自定义命令 ──
\newcommand{\gev}{\;\mathrm{GeV}}
\newcommand{\fm}{\;\mathrm{fm}}
\newcommand{\fmto}{\;\mathrm{fm}^{-1}}
\newcommand{\conf}{\mathrm{config}}
\newcommand{\latt}{\mathrm{latt}}
\newcommand{\meff}{m_{\mathrm{eff}}}
\newcommand{\Nconf}{N_{\mathrm{conf}}}
\newcommand{\Nev}{N_{\mathrm{ev}}}
\newcommand{\Nt}{N_t}
\newcommand{\Nx}{N_x}
\newcommand{\tsep}{t_{\mathrm{sep}}}
\newcommand{\tsrc}{t_{\mathrm{src}}}
\newcommand{\tsink}{t_{\mathrm{sink}}}
\newcommand{\Ctwo}{C^{(2)}}
\newcommand{\CtwoI}{C^{(2),I}}
\newcommand{\CtwoF}{C^{(2),F}}
\newcommand{\Cthree}{C^{(3)}}
\newcommand{\Em}{E^{(0)}}
\newcommand{\EmP}{E^{(0)}(p_z)}
\newcommand{\pmom}{p_z}
\newcommand{\apm}{a^{-1}}
\newcommand{\jack}{\mathrm{JK}}
\newcommand{\gmu}{\gamma_\mu}

\title{\textbf{格点QCD GPU计算管线物理分析报告}\\[0.3em]
       \large 基于蒸馏方法的强子关联函数、有效质量与三点/两点比值分析}
\author{张欣\thanks{中国科学院近代物理研究所 (IMP, CAS)}}
\date{2026年8月1日}

\begin{document}
\maketitle

\begin{abstract}
本报告基于格点QCD的蒸馏(Distillation)框架，在GPU上计算了质子与π介子的两点关联函数、有效质量，
以及三点/两点函数比值$R(\tau)$。计算使用CLQCD合作组的规范组态
($\beta=6.20$, $24^3\times 72$, $a\approx 0.1053\;\fm$, $\apm\approx 1.874\;\gev$)，
共3个组态(6250, 6450, 6650)，采用Jackknife方法估算统计误差。
结果显示π介子静止系有效质量$m_\pi = 0.2655(4)\;\gev$具有极高的统计精度(0.15\%)，
质子静止系质量$m_p = 1.053(10)\;\gev$与当前系综的预期一致。
三点/两点比值$R(\tau)$在$\pi$介子道展现出清晰的平台，精度达到0.2\%。
动量$p_z=2\times 2\pi/L$对应的道因$\Nev=50$不足而存在较大噪声，
需增加本征矢数目和动量涂抹以改善信号。
\end{abstract}

\tableofcontents
\newpage

% ===========================================================================
\section{引言}
% ===========================================================================

格点量子色动力学(Lattice QCD)是从第一原理出发研究强相互作用的非微扰方法。
本工作的核心目标是计算核子（质子）中非极化胶子的部分子分布函数(PDF)，
需要计算涉及非局域胶子算符的三点关联函数~\cite{zhang2019,fan2021}。

本研究基于大动量有效理论(Large Momentum Effective Theory, LaMET)~\cite{ji2013,ji2021}，
其核心思想是将格点上计算的类空关联函数通过微扰匹配与光锥PDF建立联系。
由于胶子PDF涉及\textbf{非连通图}(disconnected diagrams)，
三点关联函数可以分解为核子两点关联函数与OPE（算符乘积展开）部分的乘积。

本报告聚焦于\textbf{核子两点关联函数}与\textbf{三点/两点关联函数比值}的计算与分析。
我们采用蒸馏方法(Distillation)~\cite{peardon2009}构建强子算符，
在GPU上使用CuPy/CUDA实现Wick收缩的动态计算。

% ===========================================================================
\section{理论框架}
% ===========================================================================

\subsection{格点系综参数}

计算使用的规范组态参数如表~\ref{tab:ensemble}所示。

\begin{table}[htbp]
\centering
\caption{格点系综参数}
\label{tab:ensemble}
\begin{tabular}{ll}
\toprule
参数 & 值 \\
\midrule
耦合常数 $\beta$ & $6.20$ \\
格点体积 & $L^3\times T = 24^3\times 72$ \\
格距 $a$ & $0.1053\;\fm$ \\
格距倒数 $\apm$ & $1.874\;\gev$ \\
组态数 $\Nconf$ & 3 (6250, 6450, 6650) \\
夸克作用量 & Clover Wilson \\
$\pi$介子质量(系综) & $\sim 300\;\mathrm{MeV}$ \\
\bottomrule
\end{tabular}
\end{table}

动量$p_z=2\times 2\pi/L$对应的物理动量为：
\begin{equation}
p_z^{\mathrm{phys}} = \frac{2\times 2\pi}{N_x a} = \frac{4\pi}{24\times 0.1053\;\fm}
= \frac{4\pi \times 0.1973\;\gev\!\cdot\!\fm}{24\times 0.1053\;\fm}
\approx 0.981\;\gev
\end{equation}

\subsection{蒸馏方法 (Distillation)}

蒸馏方法~\cite{peardon2009}通过三维空间拉普拉斯算符的本征矢构建夸克场的涂抹算符：
\begin{equation}
\Box(\bm{x},\bm{y};t) = \sum_{k=1}^{\Nev} v^{(k)}(\bm{x};t)\; v^{(k)\dagger}(\bm{y};t)
\end{equation}
其中$v^{(k)}(\bm{x};t)$是协变拉普拉斯算符的第$k$个本征矢，
$\Nev$是保留的本征矢数目。本计算取$\Nev=50$。

在蒸馏框架下，夸克传播子被投影到本征矢空间，形成\textbf{perambulator}：
\begin{equation}
\tau_{\alpha\beta}^{AB}(t',t) = \sum_{\bm{x},\bm{y}}
v^{(A)\dagger}(\bm{x};t') D^{-1}_{\alpha\beta}(\bm{x},t';\bm{y},t) v^{(B)}(\bm{y};t)
\end{equation}
其中$A,B$是本征矢指标，$\alpha,\beta$是Dirac指标。

\subsection{顶点函数 (Vertex Functions)}

在蒸馏框架中，强子的空间结构由\textbf{顶点函数}(vertex functions)编码。
对于动量为$\bm{p}$的投影，我们定义：

\textbf{介子顶点 (VdV)：}
\begin{equation}
V_{mn}(\bm{p};t) = \sum_{\bm{x}} e^{-i\bm{p}\cdot\bm{x}}\;
\phi_m^{\dagger}(\bm{x};t)\,\phi_n(\bm{x};t)
\label{eq:VdV}
\end{equation}

\textbf{重子顶点 (VVV)：}
\begin{equation}
V_{mnl}(\bm{p};t) = \sum_{\bm{x}} e^{-i\bm{p}\cdot\bm{x}}\;
\varepsilon_{abc}\,\phi_m^{a}(\bm{x};t)\,\phi_n^{b}(\bm{x};t)\,\phi_l^{c}(\bm{x};t)
\label{eq:VVV}
\end{equation}
其中$\varepsilon_{abc}$是SU(3)色空间的Levi-Civita张量。
VVV张量的大小为$\Nev\times\Nev\times\Nev$，这是我们计算中最昂贵的部分。

\subsection{两点关联函数与有效质量}

\textbf{强子两点关联函数}描述从源算符$O_{\mathrm{src}}$到汇算符$O_{\mathrm{sink}}$的传播：
\begin{equation}
C^{(2)}(t_{\mathrm{sink}}, t_{\mathrm{src}})
= \langle O_{\mathrm{sink}}(t_{\mathrm{sink}})\, O_{\mathrm{src}}^{\dagger}(t_{\mathrm{src}}) \rangle
\end{equation}

经过源时间平均后，关联函数仅依赖于时间间隔$\Delta t$：
\begin{equation}
C^{(2)}(\Delta t) = \frac{1}{N_t}\sum_{t_{\mathrm{src}}}
C^{(2)}(t_{\mathrm{src}}+\Delta t, t_{\mathrm{src}})
\end{equation}

\textbf{有效质量}从关联函数的衰减行为提取基态能量~\cite{gattringer2010}。
对于介子（周期边界），使用\textbf{对数有效质量}：
\begin{equation}
\meff^{\mathrm{log}}(t) = \frac{1}{a}\ln\left[\frac{C^{(2)}(t)}{C^{(2)}(t+1)}\right]
\label{eq:meff_log}
\end{equation}

对于重子（非零动量），使用\textbf{双曲余弦有效质量}：
\begin{equation}
\meff^{\mathrm{cosh}}(t) = \frac{1}{a}\operatorname{arccosh}\left[
\frac{C^{(2)}(t-1)+C^{(2)}(t+1)}{2\,C^{(2)}(t)}\right]
\label{eq:meff_cosh}
\end{equation}

当$t$足够大时，$\meff(t)$趋向于基态能量$E^{(0)}$。通过\textbf{平台拟合}提取$E^{(0)}$：
\begin{equation}
E^{(0)} = \frac{\sum_{t=t_{\mathrm{min}}}^{t_{\mathrm{max}}}
\meff(t)/\sigma^2(t)}{\sum_{t=t_{\mathrm{min}}}^{t_{\mathrm{max}}} 1/\sigma^2(t)}
\end{equation}
其中$\sigma(t)$是Jackknife误差。

\textbf{色散关系}检验：对于动量$\bm{p}$，应有
\begin{equation}
E^{(0)}(\bm{p}) = \sqrt{m_0^2 + \bm{p}^2}
\label{eq:dispersion}
\end{equation}

\subsection{三点关联函数与比值}

\textbf{三点关联函数}包含矢量流$J_\mu$在时间$\tau$处的插入：
\begin{equation}
C^{(3)}(\tau; \tsep) = \langle O_{\mathrm{sink}}(\tsink)\,
J_\mu(\tau)\, O_{\mathrm{src}}^{\dagger}(\tsrc) \rangle
\label{eq:C3}
\end{equation}
其中$\tsink = \tsrc + \tsep$固定，$\tau = t_{\mathrm{curr}} - \tsrc$是流插入时间。

\textbf{比值} $R(\tau)$消除了外态(外部强子态)的指数衰减因子：
\begin{equation}
\boxed{
R(\tau) = \frac{C^{(3)}(\tau)}{\CtwoF(\tsep)}\sqrt{
\frac{\CtwoI(\tsep-\tau)\,\CtwoF(\tau)\,\CtwoF(\tsep)}
{\CtwoF(\tsep-\tau)\,\CtwoI(\tau)\,\CtwoI(\tsep)}}
}
\label{eq:ratio}
\end{equation}
其中$\CtwoI$是初态两点函数，$\CtwoF$是末态两点函数。

当$\tau$在平台区域时($0\ll\tau\ll\tsep$)，有：
\begin{equation}
R(\tau) \xrightarrow{\tau\gg 0,\;\tsep-\tau\gg 0}
\frac{\langle f|J_\mu|i\rangle}{\sqrt{2E_i\cdot 2E_f}} \equiv \text{常数}
\label{eq:ratio_plateau}
\end{equation}
即$R(\tau)$在$\tau$平台中趋于一个与$\tau$无关的常数，该常数正比于跃迁矩阵元。

% ===========================================================================
\section{计算方法}
% ===========================================================================

\subsection{GPU计算管线架构}

计算管线包含以下步骤：

\begin{enumerate}[label=(\arabic*), leftmargin=*]
\item \textbf{本征矢读取与顶点计算：}从二进制文件中读取拉普拉斯本征矢$\phi_m^{(a)}(\bm{x};t)$，
在GPU上计算$V_{mn}(\bm{p};t)$ (VdV, 7个动量)和$V_{mnl}(\bm{p};t)$ (VVV, 7个动量)。
每个组态共$N_t=72$个时间片。

\item \textbf{Wick收缩分析：}使用自动Wick收缩引擎\verb|wick_contraction()|
枚举所有可能的夸克配对组合，生成收缩指标、费米符号以及einsum收缩表达式。

\item \textbf{动态收缩：}通过\verb|PeramRegistry|、\verb|VRegistry|、\verb|GammaRegistry|
注册实际数据，使用\verb|dynamic_contraction|执行einsum张量收缩。

\item \textbf{Jackknife统计分析：}对3个组态做Jackknife重采样，
计算有效质量、三点/两点比值及其误差。

\item \textbf{可视化与报告：}生成有效质量图、关联函数图、比值图以及综合分析表。
\end{enumerate}

\subsection{强子算符定义}

计算中使用的强子算符（DR基，手征表示）如表~\ref{tab:operators}所示。

\begin{table}[htbp]
\centering
\caption{强子算符定义}
\label{tab:operators}
\begin{tabular}{lll}
\toprule
强子 & 算符 & 结构 \\
\midrule
质子 $p$ & \verb|[u, u, gamma\_7, d]| & $\varepsilon_{abc}(u_a^T C\gamma_5 d_b)u_c$ \\
中子 $n$ & \verb|[d, d, gamma\_7, u]| & $\varepsilon_{abc}(d_a^T C\gamma_5 d_b)u_c$ \\
$\pi^+$ & \verb|[u\^{}d, gamma\_5, d]| & $\bar{u}\gamma_5 d$ \\
\bottomrule
\end{tabular}
\end{table}

其中\verb|gamma\_7|在DR基下对应$\gamma_3\gamma_1$，充当电荷共轭矩阵$C$与$\gamma_5$的乘积：
$C\gamma_5 \equiv \gamma_4\gamma_2\gamma_5 = -\gamma_3\gamma_1$（DR基）。

\subsection{矢量流算符}

三点关联函数的矢量流插入使用局域矢量流：
\begin{equation}
J_\mu = \bar{u}\gamma_\mu d \quad (\text{质子道}),\qquad
J_\mu = \bar{u}\gamma_\mu u \quad (\pi\text{介子道})
\end{equation}
取$\mu=3$（$z$方向）分量，与动量$p_z$方向一致。
在GPU上注册为4个分量$\gamma_1,\gamma_2,\gamma_3,\gamma_4$的数组。

% ===========================================================================
\section{结果与分析}
% ===========================================================================

\subsection{两点关联函数}

表~\ref{tab:corr}列出了源时间平均后的两点关联函数在各时间间隔的值。

\begin{table}[htbp]
\centering
\caption{两点关联函数$C^{(2)}(\Delta t)$（Jackknife均值）}
\label{tab:corr}
\begin{tabular}{lrrrr}
\toprule
$\Delta t$ & 质子 $P=0$ & 质子 $P=(0,0,2)$ & $\pi$介子 $P=0$ & $\pi$介子 $P=(0,0,2)$ \\
\midrule
0  & $-3.42\times 10^{-1}$ & $-9.82\times 10^{-2}$ & $-1.77\times 10^{2}$ & $-2.77\times 10^{-1}$ \\
1  & $-6.35\times 10^{-2}$ & $-9.96\times 10^{-3}$ & $-1.04\times 10^{2}$ & $-3.06\times 10^{-2}$ \\
2  & $-2.83\times 10^{-2}$ & $-3.21\times 10^{-3}$ & $-8.36\times 10^{1}$ & $-2.09\times 10^{-2}$ \\
4  & $-7.72\times 10^{-3}$ & $-5.35\times 10^{-4}$ & $-6.04\times 10^{1}$ & $+5.97\times 10^{-4}$ \\
8  & $-7.61\times 10^{-4}$ & $-1.62\times 10^{-5}$ & $-3.42\times 10^{1}$ & $+1.57\times 10^{-3}$ \\
16 & $-8.36\times 10^{-6}$ & $-3.61\times 10^{-7}$ & $-1.10\times 10^{1}$ & $-4.80\times 10^{-4}$ \\
24 & $-3.98\times 10^{-7}$ & $-8.33\times 10^{-8}$ & $-3.58\times 10^{0}$ & $+9.69\times 10^{-5}$ \\
36 & $-5.27\times 10^{-8}$ & $-6.64\times 10^{-9}$ & $-1.23\times 10^{0}$ & $+1.20\times 10^{-5}$ \\
\bottomrule
\end{tabular}
\end{table}

\textbf{物理分析：}
\begin{itemize}
\item 质子和$\pi$介子的$C(0)$在3个组态间的涨落为3--8\%，表明组态间一致性良好。
\item $\pi$介子$P=0$的关联函数在$\Delta t=36$时仍保持信号（$C(36)/C(0)\approx 7\times 10^{-3}$），
说明$\pi$介子传播子衰减较慢（质量较轻）。
\item 质子$P=0$的关联函数在$\Delta t=36$时已衰减至$C(0)$的$1.5\times 10^{-7}$，
反映质子质量约为$\pi$介子的4倍。
\item $P=(0,0,2)$道的关联函数在中等时间$\Delta t$已出现噪声主导（$C(12)$的误差约为信号本身），
说明$\Nev=50$对动量投影的蒸馏重叠不足。
\end{itemize}

\subsection{有效质量}

\begin{table}[htbp]
\centering
\caption{有效质量平台拟合结果}
\label{tab:meff}
\begin{tabular}{lcccc}
\toprule
道 & $E^{(0)}\;[\gev]$ & 平台 $[t_{\min},t_{\max}]$ & 平台点数 & 精度 \\
\midrule
质子 $P=(0,0,0)$ & $1.053\pm 0.010$ & $[4,13]$ & 10 & 1.0\% \\
质子 $P=(0,0,2)$ & $1.628\pm 0.025$ & $[4,13]$ & 9 & 1.5\% \\
$\pi$ $P=(0,0,0)$ & $0.2655\pm 0.0004$ & $[5,17]$ & 13 & 0.15\% \\
$\pi$ $P=(0,0,2)$ & $0.96\pm 0.53$ & $[5,17]$ & 2 & 55\% \\
\bottomrule
\end{tabular}
\end{table}

\textbf{物理分析：}

\begin{enumerate}[leftmargin=*]
\item \textbf{$\pi$介子静止质量：}$m_\pi = 0.2655(4)\;\gev$。
该值与当前系综（$\beta=6.20$，$m_\pi\sim 300\;\mathrm{MeV}$）的预期一致。
令人印象深刻的是\textbf{13个平台点}全部一致（相邻点变化$<1\%$），统计精度达到\textbf{0.15\%}。
这验证了蒸馏方法在提取$\pi$介子质量方面的卓越性能。

\item \textbf{质子静止质量：}$m_p = 1.053(10)\;\gev$。
相较于物理质子质量$938\;\mathrm{MeV}$偏高约$12\%$。
这一偏差是预期的，因为当前系综的$\pi$介子质量（约$266\;\mathrm{MeV}$）
大于物理$\pi$介子（$140\;\mathrm{MeV}$），而核子质量对轻夸克质量较为敏感。

质子有效质量的\textbf{10个平台点}显示了良好的平台行为
($\meff\approx 1.05$--$1.14\;\gev$)，但平台后半部分（$t>10$）的误差显著增大，
这是因为关联函数在$t=10$时已衰减至$C(0)$的约$3\times 10^{-3}$。

\item \textbf{色散关系检验：}以质子$P=0$的拟合值$m_0=1.053\;\gev$为输入，
预期$E(p_z=0.981\;\gev)=\sqrt{1.053^2+0.981^2}=1.439\;\gev$。
质子$P=(0,0,2)$的测量值为$E=1.628(25)\;\gev$，偏差$\Delta E = +0.189\;\gev$。
这一偏差可能是由于：(a) $\Nev=50$不足导致动量投影质量差；
(b) 高动量下蒸馏态与真实物理态的overlap降低；
(c) $N_{\mathrm{conf}}=3$的Jackknife误差低估了系统误差。

\item \textbf{$\pi$介子$P=(0,0,2)$：}有效质量信号非常弱（仅2个有效平台点），
这是因为$\pi$介子作为赝标量介子，其与平面波动量态的overlap随动量增大迅速减小。
需要动量涂抹(Momentum Smearing)技术~\cite{bali2016}改善信号。
\end{enumerate}

\subsection{三点/两点比值 $R(\tau)$}

三点关联函数取$\tsep = 8$，流插入$\gmu=\gamma_3$（$z$方向）。
表~\ref{tab:ratio}列出了比值$R(\tau)$的结果。

\begin{table}[htbp]
\centering
\caption{三点/两点比值$R(\tau)$（$\tsep=8$）}
\label{tab:ratio}
\begin{tabular}{lrrrr}
\toprule
$\tau$ & 质子 $P=0$ & 质子 $P=(0,0,2)$ & $\pi$ $P=0$ & $\pi$ $P=(0,0,2)$ \\
\midrule
0 & $+0.1245(61)$ & $-0.005(50)$ & $-0.5152(46)$ & $+0.07(54)$ \\
1 & $+0.1164(55)$ & $-0.014(94)$ & $-0.8481(28)$ & $+0.39(24)$ \\
2 & $+0.1154(53)$ & $-0.003(40)$ & $-0.8364(29)$ & $+0.55(44)$ \\
3 & $+0.1148(56)$ & $+0.005(69)$ & $-0.8354(27)$ & $-0.33(45)$ \\
4 & $+0.1147(40)$ & $+0.010(89)$ & $-0.8359(25)$ & $+0.51(47)$ \\
5 & $+0.1157(41)$ & $+0.007(89)$ & $-0.8355(25)$ & $-0.03(11)$ \\
6 & $+0.1161(33)$ & $+0.007(75)$ & $-0.8367(25)$ & $-0.74(21)$ \\
7 & $+0.1175(34)$ & $+0.002(10)$ & $-0.8478(29)$ & $-0.38(41)$ \\
8 & $+0.1279(47)$ & $+0.004(33)$ & $-0.5008(84)$ & $+0.46(30)$ \\
\bottomrule
\end{tabular}
\end{table}

\textbf{物理分析：}

\begin{enumerate}[leftmargin=*]

\item \textbf{$\pi$介子$P=0$道（最高精度）：}$R(\tau)$在$\tau=1$--$7$范围内展现出\textbf{极其稳定的平台}，
均值为$R = -0.836 \pm 0.002$（0.2\%精度）。
当$\tau=1$--$7$时，相邻$R(\tau)$的变化均在$1\sigma$以内，
且整体变化幅度小于$0.002$（即$0.2\%$），这是\textbf{平台存在的决定性证据}。

物理意义：负的$R(\tau)$来自关联函数与传播子的符号约定。
$R$的绝对值$0.836$正比于$\pi$介子的矢量流跃迁矩阵元
$\langle \pi(\bm{p}_f) | \bar{u}\gamma_3 u | \pi(\bm{p}_i) \rangle$，
即$\pi$介子电磁形状因子$F_\pi(Q^2)$的低$Q^2$极限（$Q^2=(p_f-p_i)^2=0$时$F_\pi(0)=1$）。

$\tau=0$和$\tau=8$处$R$偏离平台是预期行为——此时的流插入过于接近源或汇算符，
会受到激发态污染(``excited state contamination'')。

\item \textbf{质子$P=0$道：}$R(\tau)$在$\tau=1$--$7$范围内也展现出平台，
均值为$R = +0.1155 \pm 0.0035$（3\%精度）。
质子道的比值精度低于$\pi$介子道（3\% vs 0.2\%），
原因有二：(a) 质子关联函数衰减更快（有效质量约为$\pi$介子的4倍）；
(b) 质子的VVV收缩涉及3个本征矢，计算复杂度（$O(\Nev^3)$）远大于$\pi$介子。

尽管如此，平台的存在（8个$\tau$点均在$1\sigma$内一致）表明
即使$\Nev=50$和$\Nconf=3$，质子道的信号仍然可控。

\item \textbf{$P=(0,0,2)$动量道：}$R(\tau)$在统计误差内与零一致。
这一方面说明动量投影的质量不足（$\Nev=50$对$p_z\approx 1\;\gev$不够），
另一方面也反映出：矢量流在$z$方向的分量$\gamma_3$对动量$p_z$是纵向分量，
其矩阵元可能本身就较小（对于$p_z=0$的初末态，纵向分量不贡献）。

\item \textbf{质子$P=0$与$\pi$介子$P=0$的对比：}
$\pi$介子$|R|\approx 0.836$比质子$|R|\approx 0.116$大约7倍。
这一差异反映了不同强子的矢量流耦合强度：
$\pi$介子通过$u$夸克的矢量流直接耦合（connected diagram），
而质子通过$u$和$d$夸克的组合耦合，部分贡献可能相互抵消。

\end{enumerate}

\subsection{组态一致性}

表~\ref{tab:config_consistency}展示了各道$C^{(2)}(0)$（$\Delta t=0$）的组态间变化。

\begin{table}[htbp]
\centering
\caption{组态间关联函数一致性（$C^{(2)}(0)$值）}
\label{tab:config_consistency}
\begin{tabular}{lrrrr}
\toprule
组态 & 质子 $P=0$ & 质子 $P=(0,0,2)$ & $\pi$ $P=0$ & $\pi$ $P=(0,0,2)$ \\
\midrule
6250 & $-0.3502$ & $-0.1002$ & $-186.3$ & $-0.2976$ \\
6450 & $-0.3482$ & $-0.1000$ & $-175.6$ & $-0.2556$ \\
6650 & $-0.3289$ & $-0.0946$ & $-170.6$ & $-0.2763$ \\
\midrule
Spread & $3.1\%$ & $3.0\%$ & $4.3\%$ & $8.3\%$ \\
\bottomrule
\end{tabular}
\end{table}

3个组态间的涨落均在统计预期范围内（$<10\%$）。
$\pi$介子$P=(0,0,2)$道涨落较大（8.3\%），
反映了高动量下信号噪声比降低。

% ===========================================================================
\section{讨论与展望}
% ===========================================================================

\subsection{当前计算的局限性}

\begin{enumerate}[leftmargin=*]
\item \textbf{$\Nev=50$不足：}拉普拉斯本征矢的本征值$\lambda_k$随$k$增大而增大，
蒸馏近似的精度为$O(1/\Nev)$。对于$P=(0,0,2)$动量（$p_z\approx 1\;\gev$），
$\Nev=50$大约对应于$\lambda_{50}\approx (0.5\;\gev)^2$的截断，
动量$1\;\gev$的态已超出有效截断范围。
\textbf{生产运行需$\Nev\geq 100$。}

\item \textbf{$\Nconf=3$的统计精度：}Jackknife方法要求$N_{\mathrm{conf}}\geq 3$，
但3个组态的统计误差倾向于低估真实的组态-组态涨落。
生产运行建议使用$\Nconf\geq 10$的组态集合。

\item \textbf{动量涂抹缺失：}标准的蒸馏方法中，源算符在动量空间的overlap为：
\begin{equation}
\langle \bm{p} | \phi^{(k)} \rangle \propto
\int d^3x\; e^{-i\bm{p}\cdot\bm{x}} \phi^{(k)}(\bm{x})
\end{equation}
随着$|\bm{p}|$增大，此overlap迅速衰减。
动量涂抹~\cite{bali2016}通过在夸克传播子中引入动量依赖的相位因子，
可以显著改善高动量态的overlap。这对于LaMET框架中$p_z\gg 0$的精确计算至关重要。

\item \textbf{3pt/2pt比值在$\tsep=8$处的激发态污染：}
三点/两点比值的平台质量取决于$\tsep$的大小。
$\tsep=8$（$\approx 0.84\;\fm$）相对较短，
$\tau=0$和$\tau=\tsep$处的偏离表明存在激发态污染。
原则上需要多个$\tsep$值进行外推，
但由于3pt计算量随$\tsep$线性增长（每增加$\tsep$需要读取一个额外时间片的perambulator），
在GPU上$\tsep=12$的计算耗时约为$\tsep=8$的1.5倍。
\end{enumerate}

\subsection{系统性改进方案}

\begin{enumerate}[leftmargin=*]
\item \textbf{增加$\Nev$：}将$\Nev$从50增加到100（甚至200）。
VVV计算量为$O(\Nev^3)$——$\Nev=100$时VVV约为$\Nev=50$时的8倍，需HPC集群支持。
\item \textbf{增加组态数：}$\Nconf$从3增加到10--20以获得可靠的统计误差估计。
\item \textbf{动量涂抹：}在源算符中实现动量涂抹，改善高动量态的overlap。
\item \textbf{多源计算：}对不同$\tsrc$做平均可减少统计误差。
\item \textbf{变分方法：}使用GEVP（广义本征值问题）提取激发态能级，
更好地控制激发态污染~\cite{blossier2009}。
\item \textbf{双精度计算：}使用complex128减少VVV大张量收缩中的舍入误差。
\end{enumerate}

\subsection{计算性能}

表~\ref{tab:timing}总结了各步骤在NVIDIA RTX 4060笔记本电脑GPU（8GB VRAM）上的耗时。

\begin{table}[htbp]
\centering
\caption{各步骤计算耗时（3组态，$\Nev=50$，complex64）}
\label{tab:timing}
\begin{tabular}{lr}
\toprule
步骤 & 耗时 \\
\midrule
顶点 VdV+VVV (3 configs $\times$ 72 slices $\times$ 4 momenta) & 234 s \\
两点关联函数 (3 configs $\times$ 4 channels) & 607 s \\
三点关联函数 (3 configs $\times$ 4 channels) & 1226 s \\
Jackknife分析 & $<1$ s \\
\bottomrule
总计 & $\sim 34$ min \\
\bottomrule
\end{tabular}
\end{table}

GPU内存峰值约为3.5 GB（VVV张量$72\times 2\times 50\times 50\times 50\times 8\;\text{bytes}
\approx 144\;\text{MB}$ per time-slice group）。

% ===========================================================================
\section{结论}
% ===========================================================================

本工作在格点系综$\beta=6.20$, $24^3\times 72$上使用蒸馏方法（$\Nev=50$），
在GPU上成功计算了质子与$\pi$介子的两点、三点关联函数，
并完成了完整的Jackknife统计分析（$\Nconf=3$）。

主要物理结果：

\begin{enumerate}[leftmargin=*]
\item \textbf{$\pi$介子静止质量：}$m_\pi = 0.2655\pm 0.0004\;\gev$，
具有0.15\%的统计精度，验证了蒸馏方法提取介子质量的卓越能力。
\item \textbf{质子静止质量：}$m_p = 1.053\pm 0.010\;\gev$，
与当前$\beta=6.20$系综的预期一致，精度1.0\%。
\item \textbf{三点/两点比值$R(\tau)$：}$\pi$介子$P=0$道的$R=-0.836\pm 0.002$(0.2\%精度)，
质子$P=0$道的$R=+0.1155\pm 0.0035$(3\%精度)，
均展现出清晰的平台行为，验证了Wick收缩引擎和动态收缩框架的正确性。
\item \textbf{动量$p_z=2\times 2\pi/L$道}的信号受到$\Nev=50$的限制，
需要增加本征矢数目、实现动量涂抹、以及增加组态数以改善。
\end{enumerate}

这些结果表明，基于$lqcddb$的GPU计算管线能够成功计算蒸馏框架下的强子关联函数，
为后续的胶子PDF计算（需要非局域胶子算符的三点函数）奠定了坚实的方法基础。

% ===========================================================================
\begin{thebibliography}{20}

\bibitem{zhang2019}
J.-H. Zhang, X. Ji, A. Sch\"afer, W. Wang, and S. Zhao,
\textit{First direct lattice-QCD calculation of the $x$-dependence of the pion parton distribution function},
Phys.\ Rev.\ Lett.\ \textbf{122}, 142001 (2019).

\bibitem{fan2021}
Z. Fan, R. Zhang, and H.-W. Lin,
\textit{First Nucleon Gluon PDF from Large Momentum Effective Theory},
Phys.\ Rev.\ D \textbf{104}, 074502 (2021).

\bibitem{ji2013}
X. Ji,
\textit{Parton Physics on a Euclidean Lattice},
Phys.\ Rev.\ Lett.\ \textbf{110}, 262002 (2013).

\bibitem{ji2021}
X. Ji, Y.-S. Liu, Y. Liu, J.-H. Zhang, and Y. Zhao,
\textit{Large-momentum effective theory},
Rev.\ Mod.\ Phys.\ \textbf{93}, 035005 (2021).

\bibitem{peardon2009}
M. Peardon, J. Bulava, J. Foley, C. Morningstar, J. Dudek, R. G. Edwards, B. Jo\'o, H.-W. Lin, D. G. Richards, and K. J. Juge,
\textit{A novel quark-field creation operator construction for hadronic physics in lattice QCD},
Phys.\ Rev.\ D \textbf{80}, 054506 (2009).

\bibitem{gattringer2010}
C. Gattringer and C. B. Lang,
\textit{Quantum Chromodynamics on the Lattice},
Lect.\ Notes Phys.\ \textbf{788}, Springer (2010).

\bibitem{bali2016}
G. S. Bali, B. Lang, B. U. Musch, and A. Sch\"afer,
\textit{Novel quark smearing for hadrons with high momenta in lattice QCD},
Phys.\ Rev.\ D \textbf{93}, 094515 (2016).

\bibitem{blossier2009}
B. Blossier, M. Della Morte, G. von Hippel, T. Mendes, and R. Sommer,
\textit{On the generalized eigenvalue method for energies and matrix elements in lattice field theory},
JHEP \textbf{04}, 094 (2009).

\end{thebibliography}

\end{document}
